% Part: model-theory
% Chapter: interpolation
% Section: definability

\documentclass[../../../include/open-logic-section]{subfiles}

\begin{document}

\olfileid{mod}{int}{def}

\olsection{परिभाष्यता प्रमेय}

अंतर्वेशन प्रमेयाचा एक महत्त्वाचा उपयोग म्हणजे बेथचे परिभाष्यता प्रमेय.
$n$-स्थानी संबंध~$R$ परिभाषित करण्यासाठी आपण $R$ चा समावेश नसलेले आणि
$n$ मुक्त !!{variable}s असलेले !!a{formula}~$!C$ देऊ शकतो. ही $!C$ च्या
साहाय्याने केलेली $R$ ची \emph{स्पष्ट} परिभाषा असेल. मग असेही म्हणता येते
की $n$-स्थानी !!{predicate}~$P$ असलेल्या !!{language} मधील
उपपत्ती~$\Sigma(P)$ मध्ये औपचारिक स्पष्ट परिभाषा समाविष्ट असेल (किंवा निदान
ती परिभाषा त्या उपपत्तीतून निष्पन्न होत असेल), म्हणजे
\[
\Sigma(P) \Entails \lforall[x_1][\dots
  \lforall[x_n][(\Atom{P}{x_1,\dots, x_n} \liff !C(x_1, \dots,
    x_n))]].
\]
तर ती उपपत्ती $P$ ला स्पष्टपणे परिभाषित करते. परंतु एखादा संबंध
परिभाषित करण्याच्या---म्हणजे तो पूर्णपणे निर्धारित करण्याच्या---मार्गांपैकी
स्पष्ट परिभाषा हा केवळ एक मार्ग आहे. कोणत्याही प्रतिमानात !!{language} च्या उर्वरित भागाचा अर्थबोध
दिला की $P$ चा अर्थबोध निश्चित होईल, अशीही एखादी उपपत्ती असू शकते.
परिभाष्यता प्रमेय असे सांगते की उपपत्ती जेव्हा या रीतीने $P$ चा अर्थबोध
निश्चित करते---जेव्हा ती $P$ ला \emph{अंतर्निहितपणे परिभाषित करते}---तेव्हा
ती त्याला स्पष्टपणेही परिभाषित करते.

\begin{defn}
समजा, !!{predicate}~$P$ चा समावेश नसलेली $\Lang{L}$ ही !!a{language} आहे. $\Lang{L}
\cup \{P\}$ मधील !!{sentence}s चा संच $\Sigma(P)$ हा $P$ ला
\emph{स्पष्टपणे परिभाषित करतो}, जर आणि केवळ जर $\Lang{L}$ मधील असे
!!a{formula}~$!C(x_1, \dots, x_n)$ असेल की
\[
\Sigma(P) \Entails \lforall[x_1][\dots
  \lforall[x_n][(\Atom{P}{x_1,\dots, x_n} \liff !C(x_1, \dots,
    x_n))]].
\]
\end{defn}

\begin{defn}
समजा, !!{predicate}s~$P$ आणि~$P'$ यांचा समावेश नसलेली
$\Lang{L}$ ही !!a{language} आहे. $\Lang{L} \cup \{P\}$ मधील
!!{sentence}s चा संच $\Sigma(P)$ हा $P$ ला \emph{अंतर्निहितपणे परिभाषित
करतो}, जर आणि केवळ जर
\[
\Sigma(P) \cup \Sigma(P') \Entails \lforall[x_1][\dots
  \lforall[x_n][(\Atom{P}{x_1,\dots, x_n} \liff \Atom{P'}{x_1,\dots,
      x_n})]],
\]
जेथे $\Sigma(P')$ हा $\Sigma(P)$ मध्ये सर्वत्र $P$ च्या जागी $P'$ घालून
मिळणारा संच आहे.
\end{defn}

दुसऱ्या शब्दांत, कोणतेही प्रतिमान $\Struct{M}$ आणि $R, R' \subseteq
\Domain{M}^n$ यांच्यासाठी, $\Expan{M}{R} \Entails \Sigma(P)$ आणि
$\Expan{M}{R'} \Entails \Sigma(P')$ हे दोन्ही खरे असतील, तर $R=R'$;
येथे $\Expan{M}{R}$ ही $\Lang{L}$ च्या $\Lang{L} \cup \{P\}$ या
विस्तारासाठीची अशी !!{structure}~$\Struct{M'}$ आहे की
$\Assign{P}{M'} = R$, आणि $\Expan{M}{R'}$ साठीही तसेच.

\begin{thm}[बेथचे परिभाष्यता प्रमेय] $\Lang{L}
  \cup\{P\}$-!!{formula}s चा संच $\Sigma(P)$ हा $P$ ला अंतर्निहितपणे
  परिभाषित करतो, जर आणि केवळ जर $\Sigma(P)$ हा $P$ ला स्पष्टपणे परिभाषित करतो.
%READERNOTE{OLFOL-043 — गोठविलेल्या इंग्रजी विधानातील “if and only if” या पदबंधात “if” गाळले आहे; मराठी विधानात आशयाला आवश्यक असलेला द्विशर्त संबंध स्पष्ट केला आहे.}
\end{thm}

\begin{proof}
$\Sigma(P)$ हा $P$ ला स्पष्टपणे परिभाषित करत असेल, तर पुढील दोन्ही निष्पन्न होतात:
\begin{align*}
  \Sigma(P) & \Entails & \lforall[x_1][\dots \lforall[x_n]
    [(\Atom{P}{x_1,\dots, x_n} \liff !C(x_1,\dots,x_n))]]\\
  \Sigma(P') & \Entails & \lforall[x_1][\dots \lforall[x_n]
    [(\Atom{P'}{x_1,\dots, x_n} \liff !C(x_1,\dots,x_n))]]
\end{align*}
आणि निष्कर्ष मिळतो. उलट दिशेसाठी, समजा $\Sigma(P)$ हा $P$ ला
अंतर्निहितपणे परिभाषित करतो. प्रथम आपण $\Lang{L}$ मध्ये !!{constant}s
$c_1$, \dots,~$c_n$ समाविष्ट करतो. मग
\[
\Sigma(P) \cup \Sigma(P') \Entails
\Atom{P}{c_1, \dots, c_n} \to  \Atom{P'}{c_1, \dots, c_n}.
\]
संहततेनुसार, असे सांत संच $\Delta_0 \subseteq \Sigma(P)$ आणि
$\Delta_1 \subseteq \Sigma(P')$ असतात की
\[
\Delta_0 \cup \Delta_1 \Entails
\Atom{P}{c_1, \dots, c_n} \to \Atom{P'}{c_1, \dots, c_n}.
\]
$!D(P)$ हा अशा सर्व !!{sentence}s $!A(P)$ चा संयोग समजू की एकतर
$!A(P) \in \Delta_0$ किंवा $!A(P') \in \Delta_1$; आणि $!D(P')$ हा अशा
सर्व !!{sentence}s $!A(P')$ चा संयोग समजू की एकतर $!A(P) \in \Delta_0$
किंवा $!A(P') \in \Delta_1$. मग $!D(P) \land
!D(P') \Entails \Atom{P}{c_1, \dots, c_n} \to \Atom{P'}{c_1, \dots, c_n}$.
%READERNOTE{OLFOL-044 — गोठविलेल्या इंग्रजी सूत्राच्या शेवटच्या अणुसूत्रात P'c_1…c_n अशी अपूर्ण कच्ची मांडणी आहे; येथे आधीपासून वापरलेले सुबद्ध अणुसूत्र पुनर्स्थापित केले आहे.}
प्रत्येक !!{predicate} $\Entails$ च्या एका बाजूला येईल अशा रीतीने याची
पुनर्मांडणी करता येते:
\[
!D(P) \land \Atom{P}{c_1, \dots, c_n} \Entails
!D(P') \to \Atom{P'}{c_1, \dots, c_n}.
\]
क्रेगच्या अंतर्वेशन प्रमेयानुसार $P$ किंवा $P'$ यांचा समावेश नसलेले असे
!!a{sentence} $!C(c_1,\dots, c_n)$ असते की
\begin{align*}
  !D(P) \land \Atom{P}{c_1, \dots, c_n} & \Entails !C(c_1,\dots, c_n); \\
  !C(c_1,\dots, c_n) & \Entails !D(P') \to \Atom{P'}{c_1, \dots, c_n}.
\end{align*}
या दोन निष्पन्नतांपैकी पहिलीवरून $!D(P) \Entails
\Atom{P}{c_1,\dots, c_n} \lif !C(c_1,\dots, c_n)$ मिळते. आणि दुसरीवरून,
$\Lang{L} \cup \{P\}$-प्रतिमान $\Expan{M}{R} \Entails !A(P)$ असेल जर
आणि केवळ जर संबंधित $\Lang{L} \cup \{P'\}$-प्रतिमान
$\Expan{M}{R} \models !A(P')$ असेल, त्यामुळे
$!C(c_1,\dots, c_n) \Entails !D(P) \lif \Atom{P}{c_1,\dots, c_n}$ मिळते,
आणि त्यावरून
\[
!D(P) \Entails !C(c_1,\dots,c_n) \to \Atom{P}{c_1,\dots, c_n}.
\]
दोन्ही एकत्र केल्यावर $!D(P) \Entails \Atom{P}{c_1,\dots, c_n}
\liff !C(c_1, \dots, c_n)$ मिळते; आणि एकस्वनिकता व सामान्यीकरण यांनुसार पुढीलही मिळते:
\[
\Sigma(P) \Entails
\lforall[x_1][\dots\lforall[x_n][(\Atom{P}{x_1,\dots, x_n} \liff
    !C(x_1,\dots, x_n))]].
\]
\end{proof}

\end{document}
